\section{Review Method}
\label{sec:method}

We conducted a structured narrative review of how contract-based and
runtime-assurance approaches represent, compose, monitor, enforce, or repair
assumptions that cross robot-software and inter-robot interfaces.  Each source
was examined with four questions:
\begin{enumerate}
  \item What is the primary object of reasoning: mission, component behavior,
    dynamical interaction, distributed state, runtime authority, or embodied
    execution?
  \item Which mechanism does the source supply: specification, refinement,
    composition, synthesis, monitoring, testing, intervention, or repair?
  \item Which kind of evidence does it report (\cref{sec:evidence-kinds}), and
    does that evidence include a physical multi-robot system?
  \item Which semantic, physical, temporal, uncertainty, evidence,
    compatibility, lifecycle, and authority obligations does it leave to
    adjacent components?
\end{enumerate}

\emph{Search.}  We searched arXiv and web-indexed IEEE, ACM, Springer, Elsevier,
and institutional repositories by title, abstract, and full text.  Core queries
combined \emph{assume--guarantee} or \emph{contract-based} with \emph{robot},
\emph{ROS}, \emph{multi-robot}, \emph{multi-agent}, \emph{mission
specification}, \emph{coordinated autonomy}, \emph{mission reliability},
\emph{probabilistic model checking}, \emph{runtime verification}, and
\emph{runtime assurance}.  Backward and forward snowballing from foundational
contract work, coordinated-mission assurance, robotic mission frameworks, ROS
contract and monitoring systems, and recent cross-layer reviews added sources
that use different terminology.

\emph{Selection.}  We included primary work that defines a contract,
compositional-assurance, runtime-monitoring, or runtime-assurance mechanism
applicable to robotics, cyber-physical systems, or interconnected agents;
reports a relevant robotics implementation; or defines a directly relevant
communication standard.  We also included surveys that organize multi-robot
modeling, communication, learning, safety, or embodied assurance.  We excluded
general design-by-contract work without an architectural or cyber-physical
connection and planning work that does not expose an interface-assurance
question.  Representative planning and middleware sources were retained only
where they clarify what an interface must carry.

\emph{Coding.}  For each source, we recorded the protected object, mechanism,
reported evidence, deployment setting, and obligations left to adjacent
components, and we grouped sources into twelve approach classes under six
protected objects (\cref{tab:taxonomy}).  The coded corpus contains 36 sources.
The remaining 15 references supply background on multi-robot coordination and
middleware, foundational contract theory, the standards used by the framework,
related surveys, and the public issue reports analyzed in
\cref{sec:middleware-evidence}.  Coverage of each obligation is coded per class
(\cref{tab:coverage}).  A class receives ``E'' when the obligation is explicit
and central in every source of the class, ``P'' when the treatment is partial,
conditional, model-internal, or present in only some sources, and ``--'' when
the obligation lies outside the stated mechanisms.  The coding describes the
cited mechanisms, not every capability of the surrounding systems.

\emph{Scope of inference.}  The search evolved through snowballing, so
PRISMA-style candidate and exclusion counts were not retained.  The review
therefore supports qualitative statements about coverage and evidence
boundaries; it supports no claim about prevalence or effect size.
